\section{近期研究（2026 年 8--11 月）}
\paragraph{DeGS：把透明度检查和重排放在硬件混合之前}
2026.08~\cite{2608_02099}。DeGS 不让每个像素处理单元在执行混合时同步检查有效片段。它先通过跨度解析找出有效任务，再重组成密集、无冲突的执行包，使后端主要用于处理有效混合，减轻稀疏覆盖和提前终止导致的阵列空转。

作者用 Tanks and Temples、Deep Blending 与 Mip-NeRF 360 场景测试从 720p 到 8K 的输出，并比较硬件架构的利用率、吞吐和能耗。相关收益对应论文中阵列、带宽和工艺配置，不等同于现有通用 GPU 上的软件端到端加速。

\paragraph{ESVR：以椭球原语直接拟合体强度并逐原语采样}
2026.08~\cite{2608_05564}。ESVR 用可微椭球拟合三维标量强度场，依结构学习与裁剪去除冗余，再沿与各原语相交的光线采样并应用传递函数。它处理的是体数据的表示和交互渲染，可与普通 3DGS 针对多视图图像的场景重建区分。

原文以大型稀疏体数据比较表示质量、渲染能力和压缩。与 TTHRESH、ZFP 和 SZ3 的对照评估数据缩减能力，作者明确将方法定位为体表示与渲染框架；该实验也不包括常规压缩的全部其他渲染版本。

\paragraph{G²ARD-GS：冻结关键几何后蒸馏紧凑高斯场}
2026.08~\cite{2608_05704}。紧凑重建除了保留当前视图质量，还可能需要给未来轨迹的外观适配和定位提供几何框架。G²ARD-GS 逐轮降低教师模型的原语预算，保护高频表面和固定锚点，再选择对受限预算有补充信息的监督视图，在不增删原语的拓扑上恢复外观。

原文用 MatrixCity block A 的 510 个留出视图评价紧凑重建，另用 495 个离轨视图评价冻结几何后的适配，用 KingsCollege 检验相机注册。五倍至三十倍压缩的结果来自这些协议。原文指出城市主评估只有一个大型场景，方法间优化预算未对齐，且只验证一种下游几何任务。

\paragraph{用显著性引导超像素合并前馈高斯}
2026.08~\cite{2608_10712}。该方法先以显著性引导的超像素切分组织逐像素原语，编码为特征高斯，再在视角之间匹配和合并。细节层级解码把输出原语数量与场景内容的复杂度关联起来。

论文分别对 DepthSplat、AnySplat 和 DA3 比较压缩比例与 DL3DV-Bench、MipNeRF360、Tanks and Temples 的新视角画质。有无输入位姿随骨干而异，具体收益还取决于压缩档位和数据集。

\paragraph{GS²CI：从单次快照压缩测量估计三维高斯}
2026.08~\cite{2608_13502}。GS²CI 用三维视觉基础模型初始化几何和相机，再令高斯渲染符合快照压缩成像的测量过程。粗阶段得到的相机轨迹在细阶段保持不变，二维基础模型从合成视角提供伪图像来校正局部外观。

实验既有从 NeRF Synthetic、DTU 等构造的合成测量，也有 SCINeRF 的真实数据。合成实验默认用 8 的压缩比和 0.25 的掩码比例。单次测量不能唯一确定所有视角的真实外观，伪视图属于模型先验提供的额外约束，需与观测重投影分开理解。

\paragraph{RoofGS：按流水线带宽与计算瓶颈协同加速渲染}
2026.08~\cite{2608_15785}。RoofGS 用 Roofline 分析区分前端的数据搬运瓶颈与光栅阶段的计算瓶颈。方法融合前端算子并压缩属性和排序键，在光栅化中同时处理两像素，并利用透明度裁剪限定后的数值范围近似指数计算。

作者在 RTX 4090 的 4K 测试中汇总 Mip-NeRF 360、Tanks and Temples 与 Deep Blending，报告平均 10.1 倍加速、616 帧每秒和 0.028 dB 的 PSNR 降低。近似指数的误差讨论取决于裁剪范围，整体速度则对应特定 GPU、图像分辨率及基线配置。

\paragraph{QuARC-GS：以锚点和量化残差压缩动态多视点视频}
2026.08~\cite{2608_18285}。QuARC-GS 以多层锚点共享初始高斯，按时间更新熵编码的量化运动，并用视图空间阈值选择需新增的原语，尽量将增量留给新出现或解除遮挡的内容。

论文在多视图视频的比较中报告每帧存储最多降至对照的约十一分之一。作者明确指出每段视频仍需高质量且成本较高的初始重建，逐帧优化也尚达不到真正实时流式交付；逐帧存储量不能直接当作传输端到端延迟。

\paragraph{极化透明度先验：从训练早期抑制冗余高斯}
2026.08~\cite{2608_22344}。该方法结合平方误差、按模糊区域增密和极化透明度正则，推动有贡献的高斯获得较明确的不透明度，同时尽早抑制低贡献基元，降低先大量增密再统一裁剪的冗余。

论文在 Mip-NeRF 360、Deep Blending 和 Tanks and Temples 比较紧凑度、训练时间与视觉质量。作者指出模糊增密使用固定 3000 至 7000 迭代的调度，复杂纹理可能丢失高频细节，原语数较低时局部渲染缺口也可能更明显。

\paragraph{KISS-GS：将原语压实与编码增益分开分析}
2026.08~\cite{2608_26948}。KISS-GS 先用 POPSpa 对训练得到的高斯做独立压实，再用 SOG-XT 编码，其中包括排序码本和协方差平滑。若编码端拥有训练视图和 GPU，还能选做考虑编码误差的微调。

作者以多个高斯预算画出率失真曲线。在 Mip-NeRF 360 的 INRIA 质量参照点，原文给出 15.7 倍、6.6 倍和额外 2.2 倍的分阶段尺寸缩减，其中最后一项需要微调。低预算下部分远背景细节消失，码本聚类和可选微调也带来编码时间。

\paragraph{非均匀量化：先估计原语重要性再合并落入同格的高斯}
2026.08~\cite{2608_28272}。该方法用原语重要性引导递归加权 Lloyd-Max 量化高斯中心，当多个原语量化到同一位置时执行重要性加权的几何合并，再编码剩余属性。量化和合并可接到点云编码器上，降低离散化和重复原语造成的编码浪费。

论文比较 G-PCC 和 V-PCC 等设置，报告测试中优于均匀体素化，并在若干配置得到更好的率失真。作者同时指出，MPEG Objects 的部分 G-PCC 配置中 Adaptive Voxelization 的率失真更好，故不能概括为所有设置都领先。

\paragraph{CC-4DGS：分别压缩动态变形和规范高斯属性}
2026.09~\cite{2609_02184}。CC-4DGS 用计算式变形场替换多分辨率哈希表，并压缩规范坐标中的高斯属性。该组合同时针对动态网络存储与点云存储，便于在保留动态渲染的前提下减少模型容量。

作者在 N3DV 和 Technicolor 与多种动态方法比较，并对 Swift4D 和本方法施加相同点云压缩以对齐存储比较。Swift4D 仍具有较高渲染速率，因此存储收益和实时速度属于需要分别衡量的量。

\paragraph{AnyGS2Mesh：从既有高斯和原分辨率视图直接生成网格}
2026.09~\cite{2609_03304}。AnyGS2Mesh 把高斯原语变成与视图对应的几何特征，以流式分块编码聚合可变数量的视图；混合深度细化将预测的局部信息与高斯渲染的尺度结合。最后经深度融合输出显式网格，无需对每个输入重新优化一套网格模型。

作者在 DTU 与 Tanks and Temples 验证，并用不同高斯重建作为输入检查可插接性。前馈免去的是逐场景的网格优化，高斯重建已存在，表面质量仍受输入表征和可用视图的证据限制。

\paragraph{分频率阶段训练多个高斯场以降低活跃原语数量}
2026.09~\cite{2609_03334}。该方法先将图像拆成低频底层和拉普拉斯残差，再从粗到细训练各频带的高斯场。训练时把不再活跃的层归档，以减少同时参与优化的原语；推断时分别渲染各层并合成结果。

论文在三类数据集的 1K 设置比较训练时间、原语量和画质。训练活跃内存的减少不能直接解释为推断只需一个场，部署仍需统计各层合成的成本。

\paragraph{从高斯到芯片：用实际负载而非原语数估计效率}
2026.09~\cite{2609_06157}。该研究将高斯存储、投影覆盖、高斯与图块关联、像素候选和有效混合分开统计，再映射到投影、排序、混合和训练更新。它用同输入 GPU 测量与文献可比性检查，展示某一阶段的计算下降可能被传输等开销抵消。

论文强调压缩和剪枝可能改变输出质量，视图一致性和更新策略也会影响重建。即便相同场景和设置也需核对计时覆盖和实际质量；更少原语、更短内核时间或更快迭代都不能单独代表端到端成果。

\paragraph{LinearMask-GS：先稳定重要性排序再按目标比例裁剪}
2026.09~\cite{2609_10095}。LinearMask-GS 把短期学习的高斯掩码激活改为线性增量，延缓过早二值化；跨视图最大池化的重要性用于在指定预算按排序裁剪，再对剩余原语细化。它将掩码学习与最后的硬选择分成两个阶段。

论文在 Mip-NeRF 360、Tanks and Temples 和 Deep Blending 检查重建并验证多种高斯骨干的迁移。解决的是掩码过早饱和与排序不稳定，真实率失真和训练耗时仍需以各场景的保留率、细化预算及相应渲染指标理解。

\paragraph{2D GauSS-MI：按面元方向估计主动采集的信息量}
2026.09~\cite{2609_21516}。该方法对增量二维高斯场中的每个面元估计视角相关的重建可靠性，以显式法向组织方向证据，再用互信息评价候选相机位置对几何与视觉重建的价值，驱动下一视点选择。

实验在 Habitat 中用带 RGB-D 的四旋翼遍历 Replica 八个房间与办公室场景。表 III 同时报告图像及几何指标，其中该方法 PSNR 为 36.87 dB，Chamfer 指标上 ActiveGS 较优。结果是仿真场景、既定传感器与导航方案下的比较。

\paragraph{CRP-GS：用锚点对应与共享先验联合压缩高斯}
2026.09~\cite{2609_23005}。CRP-GS 按特征一致性联系锚点，以对应关系建立条件编码层级，并从共享特征中聚合低冗余先验。率失真训练同时约束锚点表示与编码概率，而不是仅利用锚点距离识别可共用的信息。

作者在 Mip-NeRF 360、Tanks and Temples、Deep Blending 和 BungeeNeRF 比较存储与渲染。实现要等到第 30000 次迭代后锚点趋于稳定再开启层级训练，并继续训练 10000 次，因而结果包含特定训练计划而非任意阶段的后置压缩。

\paragraph{GARO：用梯度活动与曲率联合判断动态高斯冗余}
2026.09~\cite{2609_23509}。动态场景的增密可能积累对后续优化贡献较低的高斯，但低梯度本身不足以判断原语是否可删。GARO 将冗余判定加入自适应密度控制，先按优化活动选出候选，再用局部几何估计判断是否处在可由附近原语表达的低曲率区域。

原文在合成及真实动态场景评价速度、显存和渲染质量，包含与 D-3DGS 的 NeRF-DS 定性对照。结果支持所测数据上的速度与质量折中。几何平坦不保证外观简单，含细薄结构或高频纹理的场景仍应分别报告剪枝效果。

\paragraph{在不同动态高斯估计器的输出上独立进行压缩}
2026.09~\cite{2609_25633}。该工作把动态压缩作为对逐帧高斯序列的后处理，与高斯估计器本身分开。这样可把同一压缩器用于跟踪、半跟踪及逐帧独立的输入，分别测试不同时间对应质量下的率失真。

作者在 N3DV 六个三百帧序列比较多种输入估计器。论文的 BD-rate 以未压缩输入高斯渲染为失真参照，隔离压缩损失。因此这些 BD-rate 不能与直接对真实拍摄图像评分的端到端文献表混比。

\paragraph{观测 Gram 矩阵：用训练相机对球谐压缩定权}
2026.09~\cite{2609_28997}。该方法按训练视线方向与透明度混合贡献累积观测 Gram 矩阵，为球谐截断、码率分配和量化定权。权重反映不同系数组合在已观测方向上的颜色误差，而非简单均匀处理各阶系数。

实验用九个 Mip-NeRF 360、两个 Tanks and Temples 和两个 Deep Blending 场景的原始预训练高斯，训练相机累积统计，留出相机评价。作者指出一阶模型忽略颜色截断和共享像素中原语间的交叉项，远离捕捉视角时相对截断的优势可能消失。

\paragraph{双阶段 Morton 排序：为块式视频编码重排高斯属性}
2026.09~\cite{2609_29041}。该方法先把量化高斯位置按三维 Morton 索引排序，再用二维映射组织规则属性图。相近位置对应的属性也更易落在二维相邻块，便于传统块式编码器利用空间相关性。

论文在 GSCodec Studio 上比较 PLAS、普通 Morton 与提出的排序，采用 MPEG 测试条件中的静态单帧和动态多帧，列出率失真及计算结果。结论依赖选用的平台、码流与场景运动，不能仅据一种排列推断所有动态资产的压缩收益。

\paragraph{ADATEX4D：在冻结变形骨干上为纹理按方向分配容量}
2026.09~\cite{2609_29963}。ADATEX4D 不再给各动态高斯分配同样大小的纹理，而是依据可见归一梯度、时间峰值和变形尺度对纹理两轴分别扩容，同时冻结已有四维高斯骨干。

在 12 个场景汇总中，原文相比 Textured 4DGS 报告纹理尺寸由 109.3 MiB 降至 51.7 MiB、完整模型存储下降 24.4\%，但平均渲染由 100.8 降至 84.1 帧每秒、训练时间上升。作者指出当前稠密纹理实例化限制运行时收益，不能把省存储视为提速。

\paragraph{实用高斯压缩：在不依赖空间聚合的条件下整数化解码}
2026.09~\cite{2609_30245}。这项工作由坐标派生几何上下文，再对紧凑锚点潜变量与属性作因果编码。线性变换与激活构成的上下文网络配合率失真学习、量化感知训练和整数推理，简化解码并降低平台浮点差异。

论文在多个公开场景比较码流大小、渲染质量和总编码、解码时间，区分不同训练预算与量化感知配置。跨平台一致性是实验所用整数解码路径的结果，不能据此对未测硬件或任意码流作兼容性保证。

\paragraph{用采集规模估计容量，再按训练质量调整高斯数量}
2026.09~\cite{2609_31248}。固定一百万高斯的默认预算对不同大小的采集不一定合适。这项研究在 LiteGS 骨干上将容量上限和实际使用量拆开：初始容量按总采集像素及视图重叠决定，训练再根据拟合质量增减原语。前者反映观察了多少表面，后者反映这些表面需要多少局部细节。

原文在 13 个标准场景和 8 个较大场景使用同一配置作比较。在标准基准，bicycle 比 LiteGS 多用 18\% 高斯，stump 则少用 14\%，两者 PSNR 均有提升；truck 缩到 LiteGS 的 0.68 倍数量时 PSNR 差距约 0.1 dB，但 LPIPS 较差。紧凑度与质量的取舍依场景和指标而变，结果针对离线场景训练。

\paragraph{无线边缘联邦高斯：用重要度和时延选择端侧训练子集}
2026.09~\cite{2609_32177}。该方法在每台设备按原语重要度与开销选择可训练的高斯，以适应异构计算、存储与通信资源。不同设备剪枝后原语结构不再相同，服务器按上传的掩码恢复参数与梯度的共同索引，之后才能做结构一致的联邦聚合。

论文把真实的 Mill19 Building、Rubble 和 UrbanScene3D Polytech 图像不均匀分给模拟设备，再比较无线边缘设置下的方案。真实图像与模拟链路需要区分；该实验并未对运营网络部署、链路故障或联邦更新的隐私性作保证。

\paragraph{FeCoSplat：用初次高斯重建反馈指导第二阶段编码}
2026.09~\cite{2609_33330}。FeCoSplat 先编码多视图的几何相关特征，解码的第一阶段高斯状态用来指导发送端补发未充分表示的信息。两步共享变码率特征编码器和隐式高斯状态，接收端用轻量交互恢复可渲染原语。

论文在 RE10K 与 ACID 的率失真比较中分别考察码率和渲染质量。反馈描述的是方法内部对第一次重建的分析，不能直接解释为测过网络往返时延；发送端的第二阶段也仍需消耗计算。

\paragraph{SAGA：以粗细四维锚点压缩体积视频}
2026.09~\cite{2609_33969}。SAGA 让稀疏粗锚点提供细锚点的共享特征，再从细锚点解码高斯运动和外观。遗漏的局部细节通过锚点坐标中的哈希残差补偿。熵编码按层组织已解码的时空上下文，并用固定容量记忆减少重复信息。

论文在 Neu3D、Panoptic Sports 和 MPEG MIV 的条件下比较码率失真、解码与新视图质量。实时解码并不意味着实时编码，方法仍要逐场景优化。快速运动、拓扑变化和稀少的相机覆盖是作者指出的边界。

\paragraph{全景高斯编码：让层级球面原语的保留由率失真决定}
2026.09~\cite{2609_34367}。方法将原语位置锚定在层级 HEALPix 网格，无需为每个存在的原语另传坐标。量化透明度所导出的存在掩码与熵模型共同决定哪些原语的失真下降值得其码率，可解码整个球面或视口。

论文主要在 100 张 2048×1024 的等距柱状全景测试，并用 100 张 SUN360 图像验证跨数据集表现。评价用球面加权 PSNR 和六个视口指标，避免平面指标过度加权极区。这些是静态全景编解码实验，不能直接代表动态头显传输时延。

\paragraph{NRF-GS：用共享神经残差提高外观表达效率}
2026.09~\cite{2609_37115}。标准 3DGS 给每个原语单独存储解析球谐系数。NRF-GS 保留原语的空间表示与基本颜色，为高斯配置潜码，用共享神经残差场预测视角相关外观。方向编码再按低频与高频分开，从而把反射细节集中交由共享模型表达。

论文将参数表达能力的等内存模拟与完整重建实验分开，在 Mip-NeRF 360、DL3DV 与 Tanks and Temples 报告高斯数、训练时间及图像指标。原文所称原语最多减少 50\% 是所测设置的上限；共有网络成本也应统计。补充表标明某个 Mip-NeRF 360 对照因基线显存不足排除了 Bicycle 和 Garden。

\paragraph{TSGL：由解码端可见属性重建高斯编码图}
2026.09~\cite{2609_38635}。TSGL 先由真实属性建立教师图，再用解码端已知的输入训练小学生网络重建近似邻接，以图傅里叶基压缩属性而无需传教师图。场景级的通道变换另用于移除高阶球谐系数的跨通道相关。

论文用 Mip-NeRF 360、Tanks and Temples 与 Deep Blending 的 13 个场景测试，分别报告数据集平均尺寸和渲染指标。作者在 Mip-NeRF 360 报告 29.32 MB，而 MesonGS 为 34.41 MB。结论依据所比较的后训练方法和测试渲染设置，图估计所需的学习并未被取消。

\paragraph{ByteSplat：减少分布式高斯训练的显存和跨卡流量}
2026.10~\cite{2610_02851}。ByteSplat 融合光栅化前后向处理以减少中间值访问，用硬件感知剪枝适配共享内存，并略去跨卡消息中的零梯度，在原语归属 GPU 汇总其更新。优化目标同时覆盖片外带宽和互连。

论文在各有 8 块 GPU 的 A6000 与 A100 服务器上，针对六个数据集测量含启动时间的端到端训练，报告最高 6.1 倍加速。最高加速属于特定工作负载，不能视为每个场景或不同互连的固定收益。

\paragraph{PocketSplat：先为世界坐标潜单元分配预算再解码}
2026.10~\cite{2610_03192}。PocketSplat 将经标定的多视图证据汇合到世界坐标的局部潜单元，按目标预算精确分配整数配额，仅解码预算内的高斯。这样可以直接输出预定数量的原语，而不是生成密集高斯后临时截取。

论文在 DL3DV 用四个条件视图和 50 个评价帧评估、在 Tanks and Temples 与 MegaDepth 检查域外迁移，另对 Mip-NeRF 360 做设备侧测试。方法要求相机内参与位姿，端侧结果对应文中使用的设备和可部署对照。

\paragraph{Mobile-4DGS：把动态高斯部署约束写进表示与调度}
2026.10~\cite{2610_05289}。Mobile-4DGS 将外观修正预先烘焙进低阶球谐参数，由多视角误差指导高斯增密与裁剪。动态部分采用逐原语的显式时间运算，静态部分是零运动的特例。移动播放时只在深度排序仍满足几何条件时复用静态绘制列表。

论文在静态及动态场景测图像质量、体积、训练和实机吞吐，目标是资源受限下的取舍。缓存排序只适用于通过相机运动界限检查的帧，轻量渲染结果也不能直接推成不限资源设定下所有指标领先。

\paragraph{SteadySplats：以空间复用和方差训练降低随机渲染噪声}
2026.10~\cite{2610_05576}。SteadySplats 将低样本随机高斯渲染的噪声分为估计方差与样本数问题。一方面复用邻域和历时样本，另一方面训练时约束同一视线上贡献颜色的方差，使少量抽样也能接近按序透明度合成的结果。

论文以 7 个 Mip-NeRF 360 场景的 21 条相机轨迹、3152 帧验证，并将按序混合渲染作为参照。作者指出方差训练会稀释光度损失、略降充分收敛时的部分峰值质量指标；样本复用也并非零开销。

\paragraph{GSCV：用跨帧相容的属性图接入标准视频编码器}
2026.10~\cite{2610_07795}。GSCV 以参考布局约束相邻帧的属性排列，使同一三维结构在二维高位深属性图中尽量能被帧间预测跟踪。属性图再由视频编码器压缩，以利用已有标准而保持较高的精度上限。

论文比较 MPEG 测试中的已追踪和半追踪高斯序列，以 GSCodec Studio 与 G-PCC 为对照，报告码率和重建质量。作者指出重排后的属性图不同于自然视频，通用编码器仍会限制压缩率；多个视频属性流也会改变系统整合开销。

\paragraph{混合分层光线追踪：按区域组织高斯和三角形求交}
2026.10~\cite{2610_13498}。该工作按空间局部的密度选择加速结构，复用相邻光线的遍历状态，并在混合求交路径复用高斯和三角形的部分计算。目标是减轻局部不均衡和对邻近光线重复访问。

作者用 RTX A5000 上的 RaySplats 作为参照，并以匹配视角校验 C++ 渲染器输出。端到端速度把实测 CUDA 与模型估算的修改版 RT Core 结合，因而是架构评估，不能当成现有芯片上完整实现的实测加速。

\paragraph{SplatQ：在固定预算下联合选择相互影响的高斯}
2026.10~\cite{2610_13503}。按每个原语独立得分剪枝可能同时移除互相补偿的一组高斯，也容易让图像局部失去过多容量。SplatQ 按空间与外观划分后分配区域预算，并用颜色响应的二阶近似构造固定预算的二次选择问题，将成对影响与单点贡献共同用于删减。

原文在五个基准的 27 个场景对照重建与渲染性能，并单独衡量区域间质量不均。在其中一组消融，区域预算分配和带二次项的贪心解分别改变宏平均 PSNR。对照结果依赖预训练场景、剪枝预算和微调配置；离线求解开销也应与剪枝后的渲染帧率分别比较。

\paragraph{TETRIS-4DGS：以重要性和运动组织后训练量化}
2026.10~\cite{2610_14663}。TETRIS-4DGS 从已有动态高斯估计逐原语重要性，据此自适应量化，并按重要性与运动相关性把属性排列为二维映射，用常规图像和视频编码器分别处理相应数据。

作者从 MeetRoom 三个场景生成重要性训练标签，在 Neural 3D Video 每场景的前 300 帧评价实际压缩率、渲染失真和编解码用时。这里的后训练压缩指输入的四维高斯不用逐场景重新训练，重要性估计器仍需从已重建的资料学习。

\paragraph{PLEDGE：在声明相机域内为渐进码流提供误差界}
2026.10~\cite{2610_17605}。PLEDGE 将可接受的屏幕位置、栅格足迹、颜色和透明度偏差转换为逐高斯量化约束，并组织分层位平面，使有效截断点上的量化上界可计算。

作者将误差界验证和率失真比较分开，同时检查渐进解码。论文明确把界定义相对已剪枝模型、限制在声明的相机域，并以光栅足迹代替真实透视轮廓；对合成颜色的界还要求几何和不透明度固定。

\paragraph{RDCompGS：在高斯属性拟合时考虑二维编码率失真}
2026.10~\cite{2610_19538}。RDCompGS 在高斯表示优化中利用变换和量化系数的分布估计码率梯度，并调节率约束，使结果适合后续二维编码器，同时保留渲染质量目标。这样把原语优化与其最终能否高效编码关联起来。

论文在 Tanks and Temples、Deep Blending 和 Mip-NeRF 360 上对照 H.265/HEVC 与 JPEG XL 压缩，并比较 PSNR、SSIM 与 LPIPS。收益对应这些实际编码设置和可微码率估计，不能据此保证其他格式同幅受益。

\paragraph{前馈多视图特征编码：在接收端才构造高斯场景}
2026.10~\cite{2610_20590}。该工作先对联合外观和几何的多视图特征编码，再在接收端从解码特征前馈地产生可渲染高斯。率失真联合训练让码流控制特征信息量，并与先重建再压缩的链路区分。

作者在 RE10K 训练并在 ACID 和 DL3DV 跨域测试，前两个评估各使用 8 个输入视图，DL3DV 使用 12 个，三者均有 8 个目标视图。每个质量权重单独训练；与先生成再压缩方案的降码率不能泛化到其他输入数量或其他基线。

\paragraph{ResStreamGS：按残差敏感度分流动态视频高斯}
2026.10~\cite{2610_21121}。ResStreamGS 将视频高斯分成在一组图像内可复用的稳定部分，以及对帧间变化敏感、要逐帧编码的部分。轻量的探测适配和多属性残差评分指导选取分流比例。

论文在 Neu3D 的五段和 MeetRoom 的一段全场景视频评价，每段 300 帧、每组 20 帧，包含少量局部动态与大面积稳定背景。码流与重建比较应按作者选定的全场景序列和对照条件解释，大比例全局运动会改变可复用分量的前提。

\paragraph{LENS：在嵌入式 FPGA 协同实现定位与固定点高斯渲染}
2026.10~\cite{2610_25068}。LENS 分别实现坐标回归和有序的定点透明度混合，前者针对姿态精度，后者通过表示、递推和数据搬运协同设计减少边缘设备上的冗余带宽。

实验以 KR260 和 Jetson Orin NX 为板级比较，选用 Cambridge、Wayspots 定位和 RealEstate10K 新视图设置，训练另列为边界实验。文中的性能功耗比指相应姿态回归与渲染任务，作者不把 FPGA 定位为通用训练加速器。

\paragraph{CCT-GS：用可接入标准编码器的瓶颈图压缩高阶球谐}
2026.10~\cite{2610_26920}。CCT-GS 在多个场景训练高阶球谐图的变换，以可微编码器代理把高维属性整理为低维瓶颈图。目标场景可进一步调节前处理器，解码端在冻结主体的情况下用 LoRA 参数恢复场景相关细节。

作者按 MPEG GSC 的前向场景单帧设置、用固定留出视图比较视频式基线。训练与场景适配仍是单独步骤；兼容标准视频编码指瓶颈图的编码格式，完整高斯重建仍需论文的学习式处理。

\paragraph{KoalaGS：让同一套高斯在预算递减时保持结构}
2026.10~\cite{2610_28615}。KoalaGS 在训练中用可预先计算的图像复杂度引导原语重定位和保留排序。推理选择固定排序的前缀，从同一个模型取得不同渲染预算，而无需为每个数量重新建模。

作者以固定预算质量、完整预算区间的指标曲线以及过渡稳定性三方面比较五个基线。该方法的排序是视角无关的，并不等同于相机依赖的可见性层次结构；随机重定位仍会在重复训练之间带来一定差异。

\paragraph{ClunoGS：在同一优化中压缩高斯的语义与外观}
2026.10~\cite{2610_30013}。ClunoGS 用一个约束目标组织语义感知原语稀疏化和向量量化，以原对偶过程交替更新场景参数、压缩变量及对偶量。对语义与 RGB 同时施加监督，缓解各阶段独立压缩造成的语义丢失。

论文在 LERF 与 ScanNet-10 评价，并在同一 RGB 基础场景及相同高斯保留比例下比较先剪再量化的组合，也测试先学习语义与先压缩两种顺序。结果表明相应压缩率下的定位和分割改善，不能推出该优化对所有语言表示或码率更优。

\paragraph{BoundedGS：分别限制光栅工作区和候选梯度存储}
2026.10~\cite{2610_34233}。BoundedGS 用有负载上限的光栅单元缩短中间状态寿命，对可能可见的候选基元再做详细投影，并让反向梯度按候选紧凑存储。其设计针对高分辨率训练峰值工作区，而不是把参数和优化状态本身全部替换。

作者在固定约 600 万高斯的分辨率扫描中报告 1K 至 16K 的渲染工作区保持约 1.073 至 1.077 GiB，在五个原生 5K 场景单独报告完整训练的工作区下降。工作区按附录用分配峰值扣除模型参数与优化器等项目计算，不是整机内存或显存总预算。

\paragraph{ParetoGS：以稀疏适配锚点连接多个渲染预算}
2026.10~\cite{2610_34235}。ParetoGS 以训练完成的完整模型为起点，交替收缩原语集合并在固定成员内调整参数，形成嵌套的预算锚点。持久身份让锚点可共享成员，残差图集编码附加状态，解码时用插值覆盖锚点之间的请求预算。

论文在 Mip-NeRF 360、Tanks and Temples 与 Deep Blending 共 13 个场景对照质量和预算，使用相同留出视角。适配状态是相对完整模型保存的额外信息，需与仅存一个全模型前缀的方法按共同基线和总存储口径比较。

\paragraph{RacerGS：利用因果组上下文压缩高斯锚点残差}
2026.10~\cite{2610_34263}。RacerGS 将哈希编码投影为紧凑空间上下文，以预测的属性均值和量化步长压缩残差。FiLM 调制的循环网络再利用已解码组帮助估计后续组的熵分布。

作者在多数据集报告相对原始 3DGS 与 Scaffold-GS 的平均体积压缩，并以相应渲染指标比较。论文指出自回归上下文必须按顺序解码属性组，新场景仍需逐场景优化，压缩比不能理解成无损或实时流式解码能力。

\paragraph{MatrixSplat：将高维特征合成映射到矩阵并行单元}
2026.11~\cite{2611_15153}。MatrixSplat 以顺序无关的加权累加改写特征渲染，并将打包原语送入线程束的矩阵乘加。分批按需加载通道以减轻共享内存压力，目标是使增加语义等输出特征时不被串行混合与重复读入限制。

论文把 Mip-NeRF 360 的渲染评价与 LERF、3D-OVS 的语义结果分别比较，还报告独立内核和端到端计时。独立内核实验的加速不等于训练或完整渲染吞吐；顺序无关近似也不能当作逐像素精确的有序透明度合成。

\paragraph{GDM-GS：把沿轨迹的高斯选择重写为空间数据查询}
2026.11~\cite{2611_20546}。GDM-GS 为结构化场景的锚点建立空间索引，在物化高斯之前查询并筛除无关锚点。相邻视点可共享已解码的数据，CPU 和 GPU 按调度共同处理锚点查询，再对各视点投影混合。

论文在 BungeeNeRF 的八个场景、1600×900 分辨率对照 Scaffold-GS。轨迹窗口和相邻视图相关性是收益来源，故空间查询、共享解码带来的加速不直接代表随机相机访问或未预知位姿的交互负载。

\paragraph{选择性冻结：缓存低梯度高斯的光栅贡献以减少训练开销}
2026.11~\cite{2611_22059}。该方法逐步冻结低梯度原语，复用其对颜色和透明度合成的贡献，并重组反向光栅以兼容缓存。专用流水阵列与冻结原语缓存面向在训练中重复访问的模式。

论文分别评估冻结后的重建质量与专用架构的速度和能耗，报告最高冻结 90\% 高斯。专用吞吐与能效数字来自架构和缓存评估，不能作为购买的通用 GPU 上软件提速或流片结果引用。

\paragraph{CueSplat：为动态高斯调度稀疏精确变形}
2026.11~\cite{2611_22095}。CueSplat 固定已训练的动态高斯骨干，用历史状态预测多数原语的新参数，仅让路由器挑出的原语执行精确变形。场景预算规定精确更新规模，间歇恢复全量变形校正累积偏差。

论文用 Technicolor、N3DV 和 NeRF-DS 共 21 段序列和四种骨干比较版间吞吐及质量，基线与加速版本共用预训练检查点和渲染器。预算在场景层面校准，跨场景平均 PSNR 表现不能保证逐帧质量，短 Technicolor 序列的刷新设置也与其他数据不同。

\paragraph{Cross3DGS：在 Vulkan 中按设备可行性选择统一训练调度}
2026.11~\cite{2611_22240}。Cross3DGS 保持投影、混合、梯度和优化器的逻辑一致，仅调整内存布局与并行方式。程序根据驱动报出的硬件限制、任务规模和已知驱动问题筛除不可行方案，再按固定偏好顺序选一种调度，不在目标设备上实时搜索最快实现。

论文报告六家厂商的八个配置能完成训练，并用数值容差测试单步更新。浮点累加次序不同不会保持逐位一致，单步测试也不包含阶段间的 MCMC 原语细化，已观察到的设备与驱动规则不能直接推广到所有同厂配置。
